\chapter{Análisis clásico de sistemas lineales}

\section{Problema de valores propios generalizado}

\emph{En desarrollo.}

\section{Polos del sistema}

Desde la TF del sistema:
\begin{align*}
    H(s) &= \frac{1}{ms^2+cs+k} ~\color{red}{\frac{1/m}{1/m}}\\
    &= \frac{1/m}{s^2+2\xi\omega_ns+\omega_n^2}
\end{align*}
\begin{itemize}
    \item $\xi$: razón de amortiguamiento
    \item $\omega_n$: frecuencia natural
\end{itemize}

¿Cómo se obtienen los polos? $\Longrightarrow$ Calcular las raíces del polinomio del denominador:

\begin{align*}
    s^2+2\xi\omega_ns+\omega_n^2=0 \qquad\text{(Ecuaci\'on caracter\'istica)}
\end{align*}

\begin{align*}
     s_{1,2}=-\xi\omega_n\pm j \underbrace{\omega_n\sqrt{1-\xi^2}}_{\omega_d} \qquad \text{(Polos del sistema de 1GDL)}
     % s_{1,2}=-\xi\omega_n\pm\omega_dj ~\text{: Polos del sistema de 1GDL}
\end{align*}

¿Es estable? $\Longrightarrow$ Depende del valor de $\xi$:
\begin{itemize}
    \item $\xi>0, ~ \text{Re}(s_{1,2})<0\Longrightarrow$ Estable asintóticamente
    \item $\xi<0, ~ \text{Re}(s_{1,2})>0\Longrightarrow$ Inestable
    \item $\xi=0, ~ \text{Re}(s_{1,2})=0\Longrightarrow$ ``Estable''
\end{itemize}

Además, el sistema oscila sí y solo si $\omega_d>0\Longrightarrow \omega_n>0$: propiedad del sistema dinámico.\\
Si $\xi>0, ~ \omega_n>0$:
\begin{figure}[H]
\centering
\includegraphics[scale=0.6]{estabilidad.png}
\end{figure}

% La Función de Respuesta en Frecuencia se obtiene a través de aplicar la transformada de Fourier:
%\begin{align*}
%    H(\omega)=\mathcal{F}\{h(t)\}
%\end{align*}

Notar que el FRF se puede obtener al reemplazar $s=j\omega$ en la TF:
\begin{align*}
    H(\omega)=H(s)\bigg|_{s=j\omega}
\end{align*}

\underline{Ejemplo}: Sistema de 1GDL
\begin{align*}
    \text{EOM: }~m\ddot{u}+c\dot{u}+ku=f(t)
\end{align*}
\begin{align*}
    \Longrightarrow \text{TF: }&~ H(s)=\frac{1}{ms^2+cs+k}=\frac{1/m}{s^2+(c/m)s+k/m}\\
    \Longrightarrow \text{FRF: }&~ H(\omega)=H(s)\bigg|_{s=j\omega}=\frac{1}{(k=m\omega^2)+j\omega c}
\end{align*}
\begin{align*}
    H(\omega)=\frac{1}{(j\omega-\lambda_1)(j\omega-\lambda_1^*)}
\end{align*}
\begin{itemize}
    \item $\{\}^*$: conjugado del número complejo
\end{itemize}
\begin{align*}
    \lambda_1, \lambda_1^*=-\xi\omega_n\pm\omega_n\sqrt{1-\xi^2}~\text{: polos del sistema}
\end{align*}

\section{Función impulso (Delta Dirac)}
La función \textit{Delta de Dirac} se denota como:
\begin{align*}
    \delta(t)=   \left\{
\begin{array}{ll}
      +\infty & t=0 \\
      0 & t\neq 0 \\
\end{array} 
\right. 
\end{align*}

% Ale: O análogamente.
%\begin{align*}
%    \delta(t-\tau)=   \left\{
%\begin{array}{ll}
%      +\infty & t=\tau \\
%      0 & t\neq \tau \\
%\end{array} 
%\right. 
%\end{align*}

\begin{figure}[H]
\centering
\includegraphics[scale=0.7]{delta.png}
\end{figure}

\underline{Propiedades}:
\begin{align*}
    \int_{-\infty}^{+\infty}\delta(t)dt=1 \quad &\text{Impulso unitario}\\
    \int_{-\infty}^{+\infty}f(t)\delta(t-T)dt=f(T) \quad &\text{Sampleo}
\end{align*}

\begin{figure}[H]
\centering
\includegraphics[scale=0.7]{dirac.png}
\end{figure}

\begin{align*}
    \Rightarrow f(T) = f(t)\ast\delta(t-T)
\end{align*}

\underline{Propiedad}: Sea $f(t)$ una función continua, entonces, se puede expresar como una integral de la convolución con la función $\delta(t)$.
\begin{align*}
    f(t)=\int_{-\infty}^{+\infty}f(\tau)\delta(t-\tau)d\tau 
\end{align*}

Esta superposición nos permite componer la función $f(t)$ como una suma infinita de impulsos unitarios con la función $\delta(t)$.
\begin{figure}[h]
\centering
\includegraphics[scale=0.7]{suma_imp.png}
\end{figure}

\section{Función de respuesta a impulso unitario}
¿Cuál es la respuesta de un sistema cuando se le somete a un impulso unitario? 
\begin{figure}[H]
\centering
\includegraphics[scale=0.7]{FRI.png}
\end{figure}

\par Para sistemas estructurales ($\xi<1$), la ecuación de movimiento sometida a un impulso unitario:
\begin{align*}
    m\ddot{x}+c\dot{x}+kx=\delta(t)\\
    x(0)=0 \quad \dot{x}(0)=0
\end{align*}

La solución a esta ecuación es la Función de Respuesta Impulso (FRI) cuando el impulso se aplica en el tiempo $\tau = 0$:
\begin{align*}
    x(t)=h(t)=\frac{1}{m\omega_d}e^{-\xi\omega_nt}\sin{(\omega_dt)}, \quad t\geq0
\end{align*}

\noindent donde $\omega_n = \sqrt{\frac{k}{m}}$ es la frecuencia natural y $\omega_d = \omega_n \sqrt{1-\xi^2}$ es la frecuencia natural amortiguada. \\

Ahora, considere una excitación $u(t)$ continua. Se puede representar como un tren de impulsos:
\begin{figure}[H]
\centering
\includegraphics[scale=0.7]{FRI2.png}
\end{figure}
Notar que cada impulso genera una respuesta en el sistema $\delta(t-\tau)\longrightarrow h(t-\tau)$.

\begin{figure}[H]
\centering
\includegraphics[scale=0.7]{FRI3.png}
/\end{figure}

La respuesta del sistema cuando se somete a este tren de impulsos se obtiene por superposición: 
\begin{align*}
    y(t)=\int_0^tu(\tau)h(t-\tau)d\tau \quad \text{``Integral de Duhamel''}\Rightarrow\text{Convolución}
\end{align*}

\underline{En resumen}: La respuesta analítica de un sistema LTI es la convolución entre el FRI $h(t)$ y la excitación $u(t)$.

\begin{align*}
    \Rightarrow y(t)=u(t)\ast h(t)
\end{align*}

\section{Transformada de Laplace}
\underline{Definición}: La transformada de Laplace de $f(t)$ es igual a:
\begin{align*}
    \mathcal{L}: ~ F(s)=\int_{-\infty}^{+\infty}f(t)e^{-\xi t}dt \quad \text{Transformada Bilateral (``two-sided'')}
\end{align*}

\noindent donde $s\in\mathbb{C}$, $s=\sigma+j\omega$, $j^2=-1$, $j$ es la unidad imaginaria.\\

\underline{Definición}: La transformada inversa de Laplace es igual a:
\begin{align*}
    \mathcal{L}^{-1}: ~ f(t)=\frac{1}{2\pi j}\int_{+\sigma-j\infty}^{+\sigma+j\infty}F(s)e^{st}ds
\end{align*}

La razón por la que se opta por trabajar con la variable de Laplace $s$ en lugar de utilizar el tiempo $t$ es la simplificación para obtener la respuesta. Al emplear la variable de Laplace, la respuesta se obtiene con multiplicaciones en vez de convoluciones.

\begin{figure}[H]
\centering
\includegraphics[scale=0.7]{laplace.png}
\end{figure}

Al obtener la respuesta ya sea por el dominio del tiempo o pasando por el dominio complejo, se obtiene exactamente el mismo resultado.
\begin{align*}
    y(t)=h(t)\ast u(t) \xRightarrow[]{\mathcal{L}} Y(s)=H(s)\cdot U(s)
\end{align*}
donde:
\begin{align*}
H(s)=\mathcal{L}\{h(t)\}=\int_0^{\infty}h(t)e^{-\xi t}dt \\
Y(s)=\mathcal{L}\{y(t)\}=\int_0^{\infty}y(t)e^{-\xi t}dt \\
U(s)=\mathcal{L}\{u(t)\}=\int_0^{\infty}u(t)e^{-\xi t}dt
\end{align*}
\begin{itemize}
    \item $h(t)$: Función de Respuesta Impulso (FRI)
    \item $H(s)$: Función de transferencia (TF)
\end{itemize}

\underline{Propiedades}: 
\begin{itemize}
    \item Superposición: 
    \begin{align*}
        \mathcal{L}\{\alpha f_1(t)+\beta f_1(t)\}=\alpha\mathcal{L}\{f_1(t)\}+\beta\mathcal{L}\{f_2(t)\}
    \end{align*}
    \item Retraso (``time-delay''): 
    \begin{align*}
        f_1(t)&\xrightarrow[]{\mathcal{L}}F_1(s)\\
        f_1(t-\tau)&\xrightarrow[]{\mathcal{L}}e^{-\xi t}F_1(s)
    \end{align*}
    \item ``Time-scaling'': 
    \begin{align*}
        f(\alpha t)&\xrightarrow[]{\mathcal{L}}\frac{1}{\alpha}F(s/\alpha)
    \end{align*}
    \item Diferenciación:
    \begin{align*}
        f(t)&\xrightarrow[]{\mathcal{L}}F(s)\\
        \frac{df(t)}{dt}&\xrightarrow[]{\mathcal{L}}s^2F(s)-sf(0^+)\\
        \frac{d^2f(t)}{dt^2}&\xrightarrow[]{\mathcal{L}}s^2F(s)-sf(0^+)-\dot{f}(0^+)
    \end{align*}
    \item Integración:
    \begin{align*}
        f(t)&\xrightarrow[]{\mathcal{L}}F(s)\\
        g(t)&\xrightarrow[]{\mathcal{L}}G(s)\\
        f(t)\ast g(t)&\xrightarrow[]{\mathcal{L}}F(s)G(s)
    \end{align*}
\end{itemize}

\underline{Aparte}: 
\begin{align*}
    H(j\omega)=\mathcal{F}\{h(t)\}
\end{align*}
\begin{itemize}
    \item $H(j\omega)$: Función de Respuesta en Frecuencia (FRF)
\end{itemize}

\begin{align*}
    H(j\omega) \xleftrightarrow{s=j\omega} H(s) 
\end{align*}

\section{Función de transferencia}
Trabajar con LT permite obtener la respuesta simplemente multiplicando la señal de entrada con la función de transferencia en vez de realizar la convolución.
\begin{figure}[H]
\centering
\includegraphics[scale=0.7]{TF.png}
\end{figure}

\underline{Definición}: $H(s)\in \mathbb{C}$ (función de transferencia)
\begin{align*}
    H(s)=\frac{\prod_{i=1}^m(s-z_i)}{\prod_{j=1}^n(s-p_j)}
\end{align*}
\begin{itemize}
    \item $z_i$: polos
    \item $p_j$: ceros
\end{itemize}

\underline{Definición}: $p_j \in \mathbb{C}$ (polos) son las raíces de la ecuación característica del sistema dinámico:
\begin{align*}
    \prod_{j=1}^n(s-p_j)=0
\end{align*}

\underline{Definición}: $z_i \in \mathbb{C}$ (ceros) son las raíces del polinomio del numerador:
\begin{align*}
    \prod_{i=1}^m(s-z_i)=0
\end{align*}

\underline{Importante}: La función de transferencia recopila \textcolor{red}{toda} la información del sistema dinámico en los coeficientes de los polinomios del numerador y denominador o equivalentemente con los polos y ceros.

\section{Plano complejo}
Para visualizar las propiedades dinámicas del sistema, vamos a trabajar en el plano complejo.

\begin{figure}[H]
\centering
\includegraphics[scale=0.7]{plano_complejo.png}
\end{figure}

Los ceros $z_i$ y polos $p_j$ son puntos sobre el plano complejo
\begin{figure}[H]
\centering
\includegraphics[scale=0.7]{polos_ceros.png}
\end{figure}

\newpage
$|H(s)|$ representa una ``superficie'' sobre el plano complejo
\begin{figure}[H]
\centering
\includegraphics[scale=0.7]{magHs.png}
\end{figure}

\underline{Aparte}: Recordar que $H(s)$ corresponde a: $H(s)=\mathcal{L}\{h(t)\}$ ($h(t)$ es FRI), por ende, $h(t)=\mathcal{L}^{-1}\{H(s)\}$. Por lo tanto, los polos nos dan información del FRI.\\

Ejemplo: Si tenemos la siguiente función de transferencia $H(s)=\frac{1}{s-a}$, se puede obtener la FRI aplicando la inversa de la transformada de Laplace $h(t)=\mathcal{L}^{-1}\{H(s)\}=e^{at}$\\

\underline{Casos}:
\begin{itemize}
    \item[(i)] $a>0, ~a\in\mathbb{R}\Longrightarrow h(t)=e^{-|a|t}$
    \begin{figure}[H]
    \centering
    \includegraphics[scale=0.7]{caso1.png}
    \end{figure}
    \item[(ii)] $a=0, ~a\in\mathbb{R}\Longrightarrow h(t)=1$
    \begin{figure}[H]
    \centering
    
\includegraphics[scale=0.7]{caso2.png}
    \end{figure}
    \item[(iii)] $a<0, ~a\in\mathbb{R}\Longrightarrow h(t)=e^{|a|t}$
    \begin{figure}[H]
    \centering
    \includegraphics[scale=0.7]{caso3.png}
    \end{figure}
    \item[(iv)] $H(s)=\frac{1}{s^2+a^2}\Longrightarrow s_{1,2}=\pm \beta j \in \mathbb{C}$\\ \\
    $h(t)=\mathcal{L}^{-1}\{H(s)\}=\sin{(\beta t)}$\begin{figure}[h]
    \centering
    \includegraphics[scale=0.7]{caso4.png}
    \end{figure}
    
    \item[(v)] $a_{1,2}=\alpha\pm\beta j \in \mathbb{C}$ (polos conjugados)\\ \\
    \begin{align*}
        H(s) &= \frac{1}{(s-\alpha)^2+\beta^2} \longrightarrow \text{Raíces del denominador}\Rightarrow a_{1,2}\\
        h(t) &= \mathcal{L}^{-1}\{H(s)\} = e^{+\alpha t}\sin{(\beta t)}
    \end{align*}
    \begin{itemize}
        \item[(a)] $\alpha<0$ \\ \\
        $h(t) = e^{-|\alpha| t}\sin{(\beta t)}$
        \begin{figure}[h]
        \centering
        \includegraphics[scale=0.7]{caso5a.png}
        \end{figure}
        
        \item[(b)] $\alpha>0$\\ \\ 
        $h(t) = e^{+|\alpha| t}\sin{(\beta t)}$
        \begin{figure}[h]
        \centering
        \includegraphics[scale=0.7]{caso5b.png}
        \end{figure}
        
        \item[(c)] $\alpha=0$\\ \\
        $h(t)=\sin{(\beta t)}\Longrightarrow$ volvemos al caso (iv).
    \end{itemize}
\end{itemize}

\underline{Volviendo al plano complejo}, podemos identificar el tipo de respuesta FRI del sistema solo con información de la posición de los polos:
\begin{figure}[H]
\centering
\includegraphics[scale=0.7]{Estabilidad1.png}
\end{figure}

\section{Estabilidad \emph{BIBO}}
\begin{figure}[H]
\centering
\includegraphics[scale=0.7]{BIBO.png}
\end{figure}
Un sistema es estable en el sentido BIBO, si y sólo si el sistema posee todos sus polos a la izquierda del plano complejo. Dicho de otra forma: $p_i=\alpha_i\pm\beta_i j$
\begin{itemize}
    \item \textcolor{ForestGreen}{El sistema es estable BIBO ssi $Re(p_i)=\alpha_i<0$}
    \item \textcolor{red}{Si el sistema posee al menos un polo cuyo $Re(p_i)=\alpha_i>0$, entonces el sistema es inestable}
\end{itemize}

\section{Implementación utilizando Matlab/Simulink}

En matlab se puede trabajar con sistemas dinámicos utilizando \emph{variables tipo sistema}. Estas variables se puede crear si se cuenta con la información de alguna de las representaciones vistas anteriormente. \\

Para State-Space se requieren las matrices $\tenq{A}$, $\tenq{B}$, $\tenq{C}$, $\tenq{D} \Longrightarrow$ matrices de tipo ``double'' \\

\par Para construir un sistema desde representación espacio-estado, se utiliza el siguiente comando:
\begin{minted}[]{matlab}
>> sys = ss(As, Bs, Cs, Ds); % Variable tipo sistema 
\end{minted}

Para representar utilizando la función de transferencia se requiere conocer los coeficientes de los polinomios del numerador y denominador de la función de transferencia:

\begin{align*}
    H(s)=\frac{b_ms^m+b_{m-1}s^{m-1}+...+b_1s+b_0}{s^n+a_{n-1}s^{n-1}+...+a_1s+a_0}
\end{align*}

\par Se utiliza el siguiente comando para construir la variable tipo sistema:
\begin{minted}[]{matlab}
>> sys = tf(num,den); % Variable tipo sistema 
\end{minted}

Donde:
\begin{align*}
    \text{num} = [b_m, b_{m-1}, ... , b_1, b_0]\\
    \text{den} = [1, a_{n-1}, ... , a_1, a_0]
\end{align*}

Ejemplo de coeficientes para sistema 1GDL:
\begin{align*}
    \text{num} = [1/m]\\
    \text{den} = [1, c/m, k/m]
\end{align*}

También, se pueden generar conociendo los polos, ceros y la ganancia del sistema:
\begin{align*}
    \text{TF: }~H(s)= k\frac{\prod_{i=1}^m(s-z_i)}{\prod_{j=1}^n(s-p_j)}
\end{align*} 
\begin{align*}
    \text{zeros} &= [z_m, z_{m-1}, ..., z_2, z_1]\\
    \text{poles} &= [p_n, p_{n-1},..., p_2, p_1]\\
    \text{gain} &= k
\end{align*} 

\par Ocupando los polos, ceros y ganancia, también se puede generar una variable tipo sistema:
\begin{minted}[]{matlab}
>> sys = zpk(ceros, polos, gain); 
\end{minted}

\underline{Matlab: Control System Toolbox}\\

El toolbox \emph{Control System} de matlab permite también extraer información y simular las variables de tipo sistema.

\begin{itemize}
    \item Valores y vectores propios del sistema:
\begin{minted}[]{matlab}
>> eig(sys)
>> eig(As)
\end{minted}
\item Polos del sistema, frecuencias naturales y razones de amortiguamiento:
\begin{minted}[]{matlab}
>> damp(sys)
>> damp(As)
\end{minted}
\item Posición de ``ceros'' y ``polos'' en el plano complejo:
\begin{minted}[]{matlab}
>> pzmap(sys)
\end{minted}
\item Gráfico magnitud y fase vs frecuencia del sistema
\begin{minted}[]{matlab}
>> bode(sys)
>> bodeplot(sys)
\end{minted}
\item Ganancia del sistema
\begin{minted}[]{matlab}
>> dcgain(sys) %% DC: direct current (caso estático)
\end{minted}

\item Análisis: respuesta a impulso y a escalón
\begin{minted}[]{matlab}
>> impulse(sys)
>> step(sys)
\end{minted}
\end{itemize}




% \section{Propiedades dinámicas}

% \section{Respuesta vibraciones libres}

% \section{Respuesta vibraciones forzadas}



